%%%%%%%%%%%%%%%%%%%%%%%%%%%%%%%%%%%%%%%%%%%%%%%%%%%%%%%%%%%%%%%%%%%%%%%%%%%%%%%%%%%%%%%%%%%%%%%%%%%%%%%%%%%%%%%%%%%%%%%%%%%%%%%%%%%%%%%%
% Sabanci University thesis draft.
%%%%%%%%%%%%%%%%%%%%%%%%%%%%%%%%%%%%%%%%%%%%%%%%%%%%%%%%%%%%%%%%%%%%%%%%%%%%%%%%%%%%%%%%%%%%%%%%%%%%%%%%%%%%%%%%%%%%%%%%%%%%%%%%%%%%%%%%
\documentclass{sabanci-template}

\usepackage{amsmath,amssymb,amsthm} 
\usepackage{mathtools}
\usepackage{stmaryrd}

\usepackage{booktabs}
\usepackage{longtable}
\usepackage{array}
\usepackage{graphicx}
\graphicspath{{Figures/}}
\usepackage{textcase}
\usepackage{algorithm}
\usepackage{algpseudocode}
\usepackage[colorinlistoftodos]{todonotes}

\hypersetup{draft=false}
\DeclareUnicodeCharacter{0301}{}
\DeclareUnicodeCharacter{00C7}{\c{C}}
\DeclareUnicodeCharacter{00E7}{\c{c}}
\DeclareUnicodeCharacter{011E}{\u{G}}
\DeclareUnicodeCharacter{011F}{\u{g}}
\DeclareUnicodeCharacter{0130}{\.{I}}
\DeclareUnicodeCharacter{0131}{\i}
\DeclareUnicodeCharacter{015E}{\c{S}}
\DeclareUnicodeCharacter{015F}{\c{s}}

\let\MakeUppercase\MakeTextUppercase

\makeatletter
\@openrightfalse
\makeatother

\raggedbottom

\makeatletter
\renewcommand{\SUdedication}{%
  \ifSU@hasdedication
    \clearpage
    \null\vfill
    \begin{flushright}
      \textit{\SU@dedication}%
    \end{flushright}
    \vfill\null
  \fi
}
\renewcommand{\SUstartmain}{%
  \mainmatter
  \onehalfspacing
}
\makeatother
\newcommand{\toolname}{Coppula}
\newcommand{\Secure}{\textsf{Secure}}
\newcommand{\Insecure}{\textsf{Insecure}}

\begin{document}
%---------------------------------------------------------------------------------------------------------------------------------------
% Students should complete/check the following fields before submission.
%---------------------------------------------------------------------------------------------------------------------------------------
\degree{Master}
\thesistitle{Verification of Probing Security for Masked Circuits: A Constructive Approach and Theoretical Study}
\tezbaslik{Maskelenmiş Devreler için Sondalama Güvenliğinin Doğrulanması: Yapıcı Bir Yaklaşım ve Kuramsal Bir İnceleme}
\authorname{Kağan Gülsüm}
\programnameEN{Computer Science and Engineering}
\programnameTR{Bilgisayar Bilimi ve Mühendisliği}
\institute{Graduate School of Engineering and Natural Sciences}
\approvaldate{21-07-2026}
\advisor{Asst. Prof.}{SÜHA ORHUN MUTLUERGİL}
\jurymember{Prof.}{ERKAY SAVAŞ}
\jurymember{Asst. Prof.}{KLAUS FREIHERR VON GLEISSENTHAL}
%---------------------------------------------------------------------------------------------------------------------------------------
% ABSTRACT/OZET
%---------------------------------------------------------------------------------------------------------------------------------------
\abstractkeywordsEN{masking, probing security, formal verification, side-channel analysis, coupling}
\abstractkeywordsTR{maskeleme, sondalama güvenliği, biçimsel doğrulama, yan kanal analizi, eşlenme}

\abstractEN{Masking defends cryptographic hardware and software implementations against side-channel attacks by splitting each secret into randomised shares. The probing model of Ishai, Sahai, and Wagner reduces the security of a masked circuit to a combinatorial question: is every set of $d$ wires, called a probe, jointly independent of the secrets? This thesis studies the exact verification of that property, namely $d$-probing security, for masked arithmetic circuits over the two-element field. We first identify the computational complexity of the problem. Computing the distribution of a single probe is \#P-complete, and the corresponding security decision problem is complete for the exact-counting class C$_=$P, so verifying even a single probe sits above the polynomial hierarchy. Hence, we verify constructively and present Coppula, a $d$-probing security checker implemented in Lean 4 featuring a layered design. At its core is a rewrite engine acting as a probe oracle, which uses substitutions and a multiplicative factoring step to expose and isolate masks buried inside products. It operates alongside a subset-space traversal that avoids explicit enumeration of every $d$-subset by partitioning the space via a set-level Vandermonde identity. Ultimately, a successful discharge is interpreted as a measure-preserving coupling of the randomness space, extending the security guarantee of one certified probe to a larger safe set of wires. The checker is sound but necessarily incomplete; we formally prove its guarantees and evaluate it on masked gadgets, notably verifiying a family of cicruits where comparable tools fail.}

\abstractTR{Maskeleme, her bir sırrı rastgeleleştirilmiş paylara bölerek kriptografik donanım ve yazılım gerçekleştirimlerini yan kanal saldırılarına karşı korur. Ishai, Sahai ve Wagner'in sondalama modeli, maskelenmiş bir devrenin güvenliğini kombinatorik bir soruya indirger: \emph{sonda} olarak adlandırılan her biri $d$ kablodan oluşan kümeler, sırlardan bağımsız mıdır? Bu tez, iki elemanlı cisim üzerindeki maskelenmiş aritmetik devreler için bu özelliğin, yani $d$-sondalama güvenliğinin, tam doğrulamasını incelemektedir. Öncelikle problemin hesaplama karmaşıklığını belirliyoruz. Tek bir sondanın dağılımını hesaplamak \#P-tamdır ve karşılık gelen güvenlik karar problemi, tam-sayma sınıfı C$_=$P için tamdır; dolayısıyla tek bir sondayı doğrulamak bile polinom hiyerarşisinin üzerinde yer alır. Bu nedenle doğrulamayı yapıcı bir biçimde gerçekleştiriyor ve Lean 4'te gerçekleştirilmiş, katmanlı bir tasarıma sahip bir $d$-sondalama güvenlik denetleyicisi olan Coppula'yı sunuyoruz. Çekirdeğinde, bir sonda kehaneti olarak işlev gören bir yeniden yazma motoru bulunur; bu motor, çarpımların içine gömülü maskeleri açığa çıkarmak ve izole etmek için ikameler ve çarpanlara ayırma adımı kullanır. Bu motor, her bir $d$-alt kümesinin açıkça numaralandırılmasından kaçınan ve uzayı bir küme-düzeyi Vandermonde özdeşliği aracılığıyla bölümleyen bir alt küme-uzayı gezinmesiyle birlikte çalışır. Nihayetinde, başarıyla kapatılan bir ispat yükümlülüğü, rastgelelik uzayının ölçü-koruyan bir eşlenmesi olarak yorumlanır ve bu da tek bir sertifikalı sondanın güvenlik garantisini daha büyük, güvenli bir kablo kümesine genişletir. Denetleyici sağlamdır ancak zorunlu olarak eksiktir; garantilerini biçimsel olarak ispatlıyor ve maskelenmiş devrelerlar üzerinde, özellikle karşılaştırılabilir araçların başarısız olduğu bir devre ailesini doğrulayarak, değerlendiriyoruz.}

%---------------------------------------------------------------------------------------------------------------------------------------
% ACKNOWLEDGEMENTS/DEDICATION
%---------------------------------------------------------------------------------------------------------------------------------------
\acknowledgement{
  I would like to thank my advisor and family.
}

\dedication{
  Dedicated to my grandparents...
}

\SUfrontmatter
\glsaddall
\SUstartmain

%---------------------------------------------------------------------------------------------------------------------------------------
% MAIN MATTER
%---------------------------------------------------------------------------------------------------------------------------------------
\import{"Chapters/Chapter/"}{1_intro}
\import{"Chapters/Chapter/"}{2_preliminaries}
\import{"Chapters/Chapter/"}{3_related-work}
\import{"Chapters/Chapter/"}{4_couplings-security}
\import{"Chapters/Chapter/"}{5_design} 
\import{"Chapters/Chapter/"}{6_implementation} 
\import{"Chapters/Chapter/"}{7_evaluation} 
\import{"Chapters/Chapter/"}{8_future-work} 
\import{"Chapters/Chapter/"}{9_conclusion} 
%---------------------------------------------------------------------------------------------------------------------------------------
% BIBLIOGRAPHY
%---------------------------------------------------------------------------------------------------------------------------------------
\SUprintbibliography
%---------------------------------------------------------------------------------------------------------------------------------------
% APPENDICES
%---------------------------------------------------------------------------------------------------------------------------------------
\makeatletter
\renewcommand{\SU@appendixchapter}[1]{%
  \stepcounter{SU@appCount}%
  \refstepcounter{chapter}%
  \clearpage
  \begin{center}
    \vspace{-4.25\baselineskip}%
    {\bfseries APPENDIX\space-\space\MakeUppercase{#1}\par}%
  \end{center}%
  \addcontentsline{toc}{chapter}{APPENDIX\space-\space\MakeUppercase{#1}}%
}
\makeatother
\SUstartappendix
\import{"Chapters/Appendix/"}{appendix}
\SUafterappendix

\end{document}
